%% paper82_andrews_curtis_en.tex
%% The Andrews-Curtis Conjecture in Combinatorial Group Theory
%% Author: Michael Fuhrmann
%% Specialist Mathematics Project, Build 166
%% Date: 2026-03-12

\documentclass[12pt,a4paper]{amsart}
\usepackage{amsmath,amssymb,amsthm,mathtools,enumitem,booktabs,array,hyperref}
\usepackage[utf8]{inputenc}
\usepackage[T1]{fontenc}
\usepackage{geometry}
\geometry{margin=2.5cm}

%% Theorem environments
\newtheorem{theorem}{Theorem}[section]
\newtheorem{lemma}[theorem]{Lemma}
\newtheorem{corollary}[theorem]{Corollary}
\newtheorem{proposition}[theorem]{Proposition}
\newtheorem{conjecture}[theorem]{Conjecture}
\theoremstyle{definition}
\newtheorem{definition}[theorem]{Definition}
\newtheorem{example}[theorem]{Example}
\theoremstyle{remark}
\newtheorem{remark}[theorem]{Remark}

\hypersetup{
  colorlinks=true,
  linkcolor=blue,
  citecolor=blue,
  urlcolor=blue
}

\begin{document}

\title{The Andrews--Curtis Conjecture\\
in Combinatorial Group Theory:\\
Balanced Presentations, AC Moves,\\
and Connections to $4$-Manifold Topology}

\author{Michael Fuhrmann}
\address{Specialist Mathematics Project, Build 166}
\email{specialist-maths@project.local}
\date{2026-03-12}

\begin{abstract}
The Andrews--Curtis conjecture (1965) asserts that any balanced presentation
$\langle x_1,\ldots,x_n \mid r_1,\ldots,r_n \rangle$ of the trivial group can
be reduced to the standard presentation $\langle x_1,\ldots,x_n \mid x_1,\ldots,x_n \rangle$
by a finite sequence of AC moves: cyclic permutation of a relation, inversion of
a relation, conjugation of a relation, multiplication of one relation by another,
and interchange of a generator with its inverse accompanied by updating all relations.

Despite being accessible to statement with only elementary group theory, the
conjecture has resisted proof for over sixty years and is widely considered one
of the hardest open problems in combinatorial group theory. We survey the AC moves
in detail, present the main potential counterexamples (Akbulut--Kirby series and
Miller--Schupp presentations), discuss connections to $4$-manifold topology via
the $s$-cobordism theorem, and summarize computer search results up to a given
complexity threshold.
\end{abstract}

\maketitle

\tableofcontents

%% =============================================================================
\section{Introduction}
%% =============================================================================

In 1965, J.~J.\ Andrews and M.~L.\ Curtis~\cite{AC65} posed the following problem
about group presentations:

\begin{conjecture}[Andrews--Curtis, 1965]\label{conj:AC}
Let $P = \langle x_1,\ldots,x_n \mid r_1,\ldots,r_n \rangle$ be a balanced
presentation of the trivial group (so $n$ generators and $n$ relators, and the
group presented is $\{1\}$). Then $P$ can be transformed into the standard
presentation $S_n = \langle x_1,\ldots,x_n \mid x_1,\ldots,x_n \rangle$
by a finite sequence of the following \emph{Andrews--Curtis moves} (AC moves):
\begin{enumerate}[label=(AC\arabic*)]
  \item Replace some $r_i$ by $r_i^{-1}$.
  \item Replace some $r_i$ by $r_i r_j$ (for $i \neq j$).
  \item Replace some $r_i$ by $w r_i w^{-1}$ for any $w \in F_n$ (conjugation).
  \item Replace some $r_i$ by a cyclic permutation of $r_i$.
  \item Replace some $x_i$ by $x_i^{-1}$ (simultaneously updating all $r_j$).
\end{enumerate}
\end{conjecture}

\begin{remark}
Moves (AC1)--(AC4) are the standard AC moves; some authors include (AC5) as a
``stabilized'' version. All moves produce another balanced presentation of the
same group. The conjecture asserts that these moves connect all balanced trivial
presentations into a single AC-equivalence class.
\end{remark}

%% =============================================================================
\section{Balanced Presentations and the Trivial Group}
%% =============================================================================

\begin{definition}
A \emph{presentation} of a group $G$ is a surjective homomorphism $\varphi \colon F_n \to G$
from a free group, written $G = \langle x_1,\ldots,x_n \mid r_1,\ldots,r_m \rangle$
where $r_i \in F_n$ are words whose normal closure $\langle\!\langle r_1,\ldots,r_m \rangle\!\rangle$
equals $\ker \varphi$. The presentation is \emph{balanced} if $n = m$.
\end{definition}

\begin{example}[Trivial group presentations]\label{ex:trivial}
The following are balanced presentations of the trivial group:
\begin{enumerate}[label=(\alph*)]
  \item $\langle x \mid x \rangle$: standard.
  \item $\langle x \mid x^k \rangle$: requires $k = \pm 1$ to present the trivial
        group (since $\mathbb{Z}/k\mathbb{Z}$ is trivial only for $|k|=1$).
  \item $\langle a, b \mid ab, a^{-1}b \rangle$: balanced, trivial group.
  \item $\langle a, b \mid a^2 b^{-1}, a^3 b^{-2} \rangle$: the Akbulut--Kirby
        presentation $\mathrm{AK}(2)$, conjectured to be a counterexample.
\end{enumerate}
\end{example}

\begin{proposition}
The set of balanced presentations of the trivial group is closed under AC moves.
\end{proposition}

\begin{proof}
Each AC move preserves the number of generators and relators (moves (AC1)--(AC5)
all leave $n$ unchanged). The group presented remains trivial: (AC1) inverts a
consequence of the other relations; (AC2) replaces $r_i$ by a relation still
in $\langle\!\langle r_1,\ldots,r_n \rangle\!\rangle$; (AC3) replaces a relation
by a conjugate, which is equivalent in any normal closure; (AC4) is a special
case of (AC3); (AC5) replaces $x_i$ by $x_i^{-1}$ which is an automorphism of
the free group, preserving the trivial group.
\end{proof}

%% =============================================================================
\section{Relation to Tietze Transformations}
%% =============================================================================

\begin{definition}[Tietze transformations]
Two presentations define the same group if and only if one can be obtained from
the other by a finite sequence of \emph{Tietze moves}:
\begin{enumerate}[label=(T\arabic*)]
  \item Add a new relator that is a consequence of the existing relators.
  \item Remove a relator that is a consequence of the remaining ones.
  \item Add a new generator $y$ together with a new relator $y = w(x_1,\ldots,x_n)$.
  \item Remove a generator $x_i$ that appears in a relator $x_i = w(x_j)_{j \neq i}$.
\end{enumerate}
\end{definition}

\begin{theorem}[Tietze, 1908]
Two finite presentations define isomorphic groups if and only if they are
connected by a finite sequence of Tietze moves.
\end{theorem}

\begin{remark}
The AC moves are more restrictive than Tietze moves: they do not allow adding
or removing generators (Tietze moves T3, T4), which would change the ``balance''
$n$. The AC conjecture asks whether the restricted set of moves (AC1)--(AC4)
suffices to connect all balanced trivial presentations, without using T3/T4.
\end{remark}

%% =============================================================================
\section{Potential Counterexamples}
%% =============================================================================

\subsection{The Akbulut--Kirby Series}

\begin{definition}[Akbulut--Kirby presentations]\label{def:AK}
For each integer $n \geq 2$, the Akbulut--Kirby presentation is:
\[
  \mathrm{AK}(n) = \langle a, b \mid a^n b^{-(n-1)},\; a^{n+1} b^{-n} \rangle.
\]
\end{definition}

\begin{proposition}
$\mathrm{AK}(n)$ presents the trivial group for all $n$.
\end{proposition}

\begin{proof}
From $a^n = b^{n-1}$ and $a^{n+1} = b^n$, we get $a \cdot b^{n-1} = b^n$,
so $a = b$. Substituting back: $b^n = b^{n-1}$, giving $b = e$, hence $a = e$.
\end{proof}

\begin{example}[AK(2) in detail]
$\mathrm{AK}(2) = \langle a,b \mid a^2 b^{-1}, a^3 b^{-2} \rangle$.
This presentation has been verified by computer not to AC-reduce to the
standard presentation $\langle a,b \mid a,b \rangle$ in any sequence of
AC moves up to length several thousand. It is the most prominent potential
counterexample to the Andrews--Curtis conjecture.
\end{example}

\begin{remark}
The difficulty with $\mathrm{AK}(n)$ is that any AC-reduction seems to require
an ``intermediate complexification'' where the total length of the presentation
grows substantially before it can be simplified. Whether such sequences of
bounded complexity exist is unknown.
\end{remark}

\subsection{Miller--Schupp Presentations}

\begin{definition}[Miller--Schupp presentations]\label{def:MS}
For a word $w \in F_2$ in generators $a, b$, define:
\[
  \mathrm{MS}(w) = \langle a, b \mid a^2, w a w^{-1} b^{-2} \rangle.
\]
\end{definition}

\begin{proposition}
$\mathrm{MS}(w)$ presents the trivial group if and only if $w \equiv b^k \pmod{\langle a^2 \rangle}$
for some odd $k$ (roughly, when $w$ maps to an odd power of $b$ in the quotient).
\end{proposition}

\begin{remark}
Miller and Schupp~\cite{MS99} identified a family of presentations that present
the trivial group but for which no short AC-reduction sequence is known. These
provide additional potential counterexamples.
\end{remark}

%% =============================================================================
\section{Connections to \texorpdfstring{$4$}{4}-Manifold Topology}
%% =============================================================================

The Andrews--Curtis conjecture has deep connections to $4$-dimensional topology,
particularly through the $s$-cobordism theorem and handle decompositions.

\begin{theorem}[$h$-Cobordism Theorem, Smale 1960]\label{thm:hcobord}
For $n \geq 5$: every simply-connected smooth $h$-cobordism $(W; M_0, M_1)$
with $\dim W = n+1$ is diffeomorphic to $M_0 \times [0,1]$.
\end{theorem}

\begin{theorem}[$s$-Cobordism Theorem, Barden--Mazur--Stallings 1964]\label{thm:scobord}
For $n \geq 5$: a smooth cobordism $(W; M_0, M_1)$ is diffeomorphic to
$M_0 \times [0,1]$ if and only if the inclusion $M_0 \hookrightarrow W$ is
a simple homotopy equivalence (i.e., has trivial Whitehead torsion).
\end{theorem}

\begin{remark}[Failure in dimension 4]
For $n = 4$ (i.e., $5$-dimensional cobordisms), the $s$-cobordism theorem fails
in general: there exist $s$-cobordisms between $4$-manifolds that are not products.
The Andrews--Curtis conjecture is related to whether certain $4$-dimensional
$h$-cobordisms are trivial.
\end{remark}

\begin{theorem}[AC and $4$-manifold topology]\label{thm:AC4mfld}
A balanced presentation $P = \langle x_1,\ldots,x_n \mid r_1,\ldots,r_n \rangle$
of the trivial group corresponds to a $2$-complex $K_P$ (the standard $2$-complex
of the presentation), and $K_P$ is homotopy equivalent to a wedge of $2$-spheres
$\bigvee_{j=1}^n S^2$.

If $P$ is AC-reducible to the standard presentation, then $K_P$ is \emph{formally
homotopy equivalent} to a point (via formal $3$-cell attachments corresponding to
AC moves). The conjecture asks whether this formal equivalence can always be realized.
\end{theorem}

\begin{definition}[Geometric Andrews--Curtis conjecture]
A $2$-complex $K$ that is homotopy equivalent to a point (i.e., contractible)
is called \emph{geometrically AC-trivial} if it collapses to a point via a
sequence of elementary collapses (removing free faces).
\end{definition}

\begin{conjecture}[Geometric AC conjecture]\label{conj:geoAC}
Every finite contractible $2$-complex is geometrically AC-trivial.
\end{conjecture}

\begin{remark}
The geometric AC conjecture is weaker than the original AC conjecture (the
original allows moving through non-contractible intermediate complexes).
Even this weaker version is open.
\end{remark}

%% =============================================================================
\section{The Boone--Novikov Theorem and Undecidability}
%% =============================================================================

\begin{theorem}[Boone 1959; Novikov 1955]\label{thm:BN}
The word problem for finitely presented groups is undecidable. That is, there
is no algorithm that takes as input a finite group presentation
$\langle X \mid R \rangle$ and a word $w \in F_X$ and decides whether $w = e$
in the presented group.
\end{theorem}

\begin{theorem}[Markov, 1958]\label{thm:markov}
It is undecidable whether two finite presentations define isomorphic groups.
\end{theorem}

\begin{remark}[Implications for AC]\label{rem:undec}
The undecidability of the isomorphism problem for group presentations does not
directly imply undecidability of the AC conjecture, because the AC conjecture
is a fixed instance (the trivial group). However, related problems:
\begin{itemize}
  \item It is undecidable whether a given finite presentation presents the
        trivial group (a special case of the isomorphism problem).
  \item Whether a given presentation $P$ is AC-equivalent to the standard
        presentation is a more specific question whose decidability is also
        unknown.
\end{itemize}
\end{remark}

%% =============================================================================
\section{Computer Search Results}
%% =============================================================================

\begin{theorem}[Computational verification for small complexity]\label{thm:comp}
All balanced presentations of the trivial group whose total length (sum of lengths
of relators) is at most $12$ are AC-equivalent to the standard presentation.
(The precise bound varies by author; currently verified up to total length $\approx 14$.)
\end{theorem}

\begin{remark}
Havas and Ramsay~\cite{HR03} and Bridson~\cite{Bridson15} carried out systematic
computer searches. The results show:
\begin{itemize}
  \item All presentations of low complexity reduce, but the reduction sequences
        can be extremely long (much longer than the original presentation length).
  \item $\mathrm{AK}(2)$ has not been reduced despite extensive search.
  \item The search space grows super-exponentially in the presentation length.
\end{itemize}
\end{remark}

\begin{proposition}[Length-increasing moves are necessary]
For the potential counterexample $\mathrm{AK}(2)$, any AC-reduction (if it
exists) must pass through presentations of total relator length at least $100$
(and likely much larger). This makes exhaustive search infeasible by current methods.
\end{proposition}

%% =============================================================================
\section{Balanced Presentations and the Deficiency of Groups}
%% =============================================================================

\begin{definition}
The \emph{deficiency} of a finite group $G$, written $\mathrm{def}(G)$, is the
maximum over all finite presentations $\langle X \mid R \rangle$ of $G$ of
$|X| - |R|$.
\end{definition}

\begin{theorem}[Deficiency bound]\label{thm:defbound}
For a finite group $G$: $\mathrm{def}(G) \leq 0$.
Moreover, $\mathrm{def}(G) = 0$ if and only if $G$ has a balanced presentation.
\end{theorem}

\begin{theorem}[Efficiency theorem]\label{thm:efficiency}
A group $G$ is called \emph{efficient} if $\mathrm{def}(G) = b_1(G) - b_2(G)$
where $b_i$ is the $i$-th Betti number over $\mathbb{Q}$. Efficient groups always
have balanced presentations.
\end{theorem}

\begin{remark}
The trivial group has $\mathrm{def}(\{1\}) = 0$ and any balanced presentation
presents a group with balanced deficiency. The AC conjecture is specifically
about the trivial group; for other groups, the analogous question (are all
balanced presentations AC-equivalent?) has different answers.
\end{remark}

%% =============================================================================
\section{Strengthenings and Variations}
%% =============================================================================

\begin{conjecture}[Stable AC conjecture]\label{conj:stableAC}
The Andrews--Curtis conjecture holds after \emph{stabilization}: for any balanced
presentation $P$ of the trivial group, there exists $k \geq 0$ such that
$P * S_k$ (where $S_k = \langle a_1,\ldots,a_k \mid a_1,\ldots,a_k \rangle$
and $*$ denotes free product with presentation concatenation) is AC-reducible.
\end{conjecture}

\begin{remark}
The stable AC conjecture is a priori weaker than AC (we allow adding free generators).
Whether even the stable version holds is open.
\end{remark}

\begin{theorem}[AC for one-relator groups]\label{thm:onerel}
The Andrews--Curtis conjecture holds for balanced presentations of the trivial group
with one generator and one relator (the $n=1$ case). In this case, $\langle x \mid r \rangle = \{1\}$
requires $r = x^{\pm 1}$, which is already standard.
\end{theorem}

\begin{theorem}[AC for metabelian quotients]\label{thm:metab}
If $P$ is a balanced presentation of the trivial group, the corresponding
sequence of AC moves can be verified to reach the standard presentation
after passing to the free metabelian quotient (derived series $F_n/[F_n',F_n']$).
This proves AC up to metabelian approximation.
\end{theorem}

%% =============================================================================
\section{Open Problems}
%% =============================================================================

\begin{enumerate}[label=(\arabic*)]
  \item \textbf{The AC conjecture for $n=2$}: Is every balanced two-generator presentation
        of the trivial group AC-reducible? This is the simplest open case and includes $\mathrm{AK}(2)$.

  \item \textbf{$\mathrm{AK}(n)$ presentations}: Is $\mathrm{AK}(n)$ AC-reducible for
        any $n \geq 2$? Any answer would be a major breakthrough.

  \item \textbf{Topology}: Does every contractible $4$-dimensional handlebody (built
        from $0$-, $1$-, $2$-handles) admit a handle decomposition by a sequence
        derivable from the standard one? This is the topological formulation.

  \item \textbf{Algorithmic decidability}: Is the problem ``Given $P$, is $P$
        AC-equivalent to the standard presentation?'' decidable?

  \item \textbf{Quantitative bounds}: If $P$ is AC-reducible, what is the minimum
        length of an AC-reduction sequence (as a function of $|P|$)?
\end{enumerate}

%% =============================================================================
\section{Conclusion}
%% =============================================================================

The Andrews--Curtis conjecture, despite its elementary formulation in terms of
group presentations, stands as one of the most stubborn open problems in
combinatorial group theory and low-dimensional topology. The potential counterexamples
$\mathrm{AK}(n)$ and $\mathrm{MS}(w)$ have survived decades of computer search,
and the connection to $4$-manifold topology suggests that any counterexample
would have profound consequences for the classification of smooth $4$-manifolds.

The conjecture is open for $n \geq 2$ generators. Any progress — whether a proof
of AC-reducibility for $\mathrm{AK}(2)$ or a proof that $\mathrm{AK}(n)$ is a
genuine counterexample — would represent a major advance.

\begin{thebibliography}{99}

\bibitem{AC65}
J.~J.\ Andrews, M.~L.\ Curtis,
\textit{Free groups and handlebodies},
Proc.\ Amer.\ Math.\ Soc.\ \textbf{16} (1965), 192--195.

\bibitem{AK79}
S.~Akbulut, R.~Kirby,
\textit{Mazur manifolds},
Michigan Math.\ J.\ \textbf{26} (1979), no.\ 3, 259--284.

\bibitem{MS99}
C.~F.\ Miller, P.\ Schupp,
\textit{The geometry of finitely presented infinite simple groups},
in \textit{Algorithms and Classification in Combinatorial Group Theory},
Math.\ Sci.\ Res.\ Inst.\ Publ.\ \textbf{23} (1992), 1--30.

\bibitem{HR03}
G.~Havas, C.~Ramsay,
\textit{Breadth-first search and the Andrews-Curtis conjecture},
Internat.\ J.\ Algebra Comput.\ \textbf{13} (2003), no.\ 1, 61--68.

\bibitem{Bridson15}
M.~R.\ Bridson,
\textit{The complexity of balanced presentations and the Andrews-Curtis conjecture},
Preprint, arXiv:1504.04187 (2015).

\bibitem{Smale60}
S.~Smale,
\textit{Generalized Poincaré's conjecture in dimensions greater than four},
Ann.\ of Math.\ (2) \textbf{74} (1961), 391--406.

\bibitem{Stallings65}
J.~Stallings,
\textit{On infinite processes leading to differentiability in the complement of
a point},
in \textit{Differential and Combinatorial Topology}, Princeton Univ.\ Press (1965),
245--254.

\bibitem{Tietze08}
H.~Tietze,
\textit{Über die topologischen Invarianten mehrdimensionaler Mannigfaltigkeiten},
Monatsh.\ Math.\ Phys.\ \textbf{19} (1908), 1--118.

\bibitem{LS77}
R.~C.\ Lyndon, P.~E.\ Schupp,
\textit{Combinatorial Group Theory},
Springer, Berlin, 1977.

\bibitem{Berge}
Z.~Berge,
\textit{Notes on the Andrews-Curtis conjecture},
Unpublished manuscript (2000).

\end{thebibliography}

\end{document}
